\section{WC 3/2/25}

\subsection*{Overview:}
This week I brought my document editor to a state of being able to create, save and recall detailed multimedia notes.

\subsection*{Achievements, Thoughts and Ideas:}
Building on last week's progress I added a number of 'extensions' to my document editor which further enhance the functionality a user has access to. I have added some premade (by other developers) plugin such as YouTube support which allows the users to seamlessly integrate a YouTube video to their notes but also some more advanced custom functionality such as my @key tag adding and / command search (which is heavily inspired by Notion as identified in my primary research, but my own implementation). 

While adding these commands I came up against a small yet potentially larger decision to the design philosophy of my software. Previously I mentioned adding tags using the @ symbol (as can be seen in the prototype from this week). Initially I had wanted this to be done using the \# symbol however this icon is usually reserved for adding headings by the markdown standard. Whilst some apps (such as obsidian) use the same key for both functions I believe this could be unclear to users. This combined with Apple favouring \# being accessed with option + 3 (a fairly unnatural key binding) I settled on the @ key for its intuitive context (you are @atting a tag to a note), easy key binding and lack of overlap.  Whilst in this case I have elected to follow preset standards I might diverge away from this in the future should I feel these become restrictive, it is definitely a fine balance between supporting a users muscle memory and making the software as polished as possible. One plugin I had hoped to add this week was an ability for users to easily type mathematical and chemical equations. This feature had been identified by a couple of respondents in my survey and I had thought would be easy to develop. I found a library called Katex which makes it easier to display these equations on the stage however I faced some challenges integrating this with TipTap and faced a number of bugs that resulted in the webapp crashing seemingly at random. Whilst this is definitely a feature I want to add it is still marked as an 'auxiliary' feature in my plans, and so I am going to postpone adding it until other core features are developed.

At this point I have a strong foundation for my 'Document' note-taking system including a wide range of plugins and features that will allow users to create rich, multimedia notes. Please find my current prototype here: \nameref{artifact:DocumentEditor}.

Throughout my research the term PWA (progressive web app) was mentioned a number of times and I spent some time researching (Please find here: \nameref{research:Webapps}) the ways I could build my webapp into one of these (making it more accessible for mobile users). While this may be out of the scope of this project I still found the research valuable and will be applicable for other projects I may be working on. Additionally, I had a look at a number of software packages that enable developers of mobile apps such as mine to package them natively for mobile users. This has a number of benefits such as allowing access to device APIs such as haptic feedback (vibrations signalling actions), camera controls and the ability to 'share' content to the app whilst maintaining the benefits of a webapp such as easy development and widespread support. That last share feature was especially interesting as I look to minimise the effort required by a user to record notes, websites, photos etc and this 2 click share wouldn't even require them to open my app saving time. Whilst outside the scope of my project for now, should have time at the end this could definitely be added.

In lesson time this week I reviewed my timeline I had previously made to progress check where I am up to. So far I am ahead of schedule and hope to keep this momentum going. To reflect my progress I have updated my timeline to reflect the current state of my project and what I hope to achieve in the coming weeks. Please find my Timeline 1 reflection and Timeline 2 here: \nameref{planning:Timeline2}


\subsection*{Power of Three:}
\textbf{This week's Tasks laid out per day}

\begin{tabularx}{\textwidth}{
  >{\raggedright\arraybackslash}X  % First column expands
  >{\centering\arraybackslash}p{1.5cm}  % Middle column has fixed width
  >{\raggedright\arraybackslash}X  % Third column expands
}
\toprule
\textbf{Goal} & \textbf{Achieved?} & \textbf{Notes} \\
\midrule


\rowcolor{gray!10}
\multicolumn{3}{l}{\textbf{Monday}} \\
Add in Tiptap + community plugins for quick feature addition & \ding{51} & \\



\rowcolor{gray!10}
\multicolumn{3}{l}{\textbf{Tuesday}} \\
Build @ tag adding system & \ding{51} & \\



\rowcolor{gray!10}
\multicolumn{3}{l}{\textbf{Thursday}} \\
Build / command searching & \ding{51} & \\


\rowcolor{gray!10}
\multicolumn{3}{l}{\textbf{Friday}} \\
Review timeline / long term plan & \ding{51} & \\

\bottomrule
\end{tabularx}


\subsection*{Plans for next week:}
With my document editor complete (for now) I want to begin building my writing focused whiteboard tool. The goal is to have it complete by the end of March so I can begin user testing and refining both tools. I will also continue to refine my document editor, adding in more plugins and features as I go.
